% TerraPulse AI — Ecuaciones del radar usadas o mencionadas en el proyecto
% Compilar con: pdflatex ecuaciones_radar.tex  (o latexmk -pdf)
\documentclass[11pt,a4paper]{article}
% Preámbulo común de la documentación LaTeX de TerraPulse AI.
% Se incluye con % Preámbulo común de la documentación LaTeX de TerraPulse AI.
% Se incluye con % Preámbulo común de la documentación LaTeX de TerraPulse AI.
% Se incluye con \input{preambulo} desde cada documento. Compilar con pdflatex (o latexmk -pdf).
\usepackage[utf8]{inputenc}
\usepackage[T1]{fontenc}
\usepackage{lmodern}
\usepackage[spanish,es-noshorthands,es-tabla,es-nodecimaldot]{babel}
\usepackage{amsmath,amssymb}
\usepackage[margin=2.3cm]{geometry}
\usepackage{array,booktabs,longtable}
\usepackage[normalem]{ulem}
\usepackage{xcolor}
\usepackage{enumitem}
\usepackage{listings}
\usepackage[hidelinks]{hyperref}

\setlength{\parindent}{0pt}
\setlength{\parskip}{0.5em}
\setlist{itemsep=2pt,topsep=2pt}
\renewcommand{\arraystretch}{1.25}
\Urlmuskip=0mu plus 1mu

% Símbolos frecuentes
\newcommand{\sz}{\ensuremath{\sigma^0}}
\newcommand{\gz}{\ensuremath{\gamma^0}}
\newcommand{\bz}{\ensuremath{\beta^0}}
\newcommand{\dsz}{\ensuremath{\Delta\sigma^0}}
\newcommand{\ok}{\ensuremath{\checkmark}}
\newcommand{\flecha}{\ensuremath{\rightarrow}}
\newcommand{\aprox}{\ensuremath{\approx}}
\newcommand{\por}{\ensuremath{\times}}

% Recuadro de nota (equivale a los ">" del Markdown)
\newenvironment{nota}
  {\begin{quote}\small\itshape}
  {\end{quote}}

% Bloques de código y comandos
\lstset{
  basicstyle=\ttfamily\small,
  breaklines=true,
  columns=fullflexible,
  keepspaces=true,
  frame=single,
  rulecolor=\color{black!30},
  backgroundcolor=\color{black!3},
  xleftmargin=4pt, xrightmargin=4pt,
  extendedchars=true,
  literate={á}{{\'a}}1 {é}{{\'e}}1 {í}{{\'i}}1 {ó}{{\'o}}1 {ú}{{\'u}}1
           {Á}{{\'A}}1 {É}{{\'E}}1 {Í}{{\'I}}1 {Ó}{{\'O}}1 {Ú}{{\'U}}1
           {ñ}{{\~n}}1 {Ñ}{{\~N}}1
}
 desde cada documento. Compilar con pdflatex (o latexmk -pdf).
\usepackage[utf8]{inputenc}
\usepackage[T1]{fontenc}
\usepackage{lmodern}
\usepackage[spanish,es-noshorthands,es-tabla,es-nodecimaldot]{babel}
\usepackage{amsmath,amssymb}
\usepackage[margin=2.3cm]{geometry}
\usepackage{array,booktabs,longtable}
\usepackage[normalem]{ulem}
\usepackage{xcolor}
\usepackage{enumitem}
\usepackage{listings}
\usepackage[hidelinks]{hyperref}

\setlength{\parindent}{0pt}
\setlength{\parskip}{0.5em}
\setlist{itemsep=2pt,topsep=2pt}
\renewcommand{\arraystretch}{1.25}
\Urlmuskip=0mu plus 1mu

% Símbolos frecuentes
\newcommand{\sz}{\ensuremath{\sigma^0}}
\newcommand{\gz}{\ensuremath{\gamma^0}}
\newcommand{\bz}{\ensuremath{\beta^0}}
\newcommand{\dsz}{\ensuremath{\Delta\sigma^0}}
\newcommand{\ok}{\ensuremath{\checkmark}}
\newcommand{\flecha}{\ensuremath{\rightarrow}}
\newcommand{\aprox}{\ensuremath{\approx}}
\newcommand{\por}{\ensuremath{\times}}

% Recuadro de nota (equivale a los ">" del Markdown)
\newenvironment{nota}
  {\begin{quote}\small\itshape}
  {\end{quote}}

% Bloques de código y comandos
\lstset{
  basicstyle=\ttfamily\small,
  breaklines=true,
  columns=fullflexible,
  keepspaces=true,
  frame=single,
  rulecolor=\color{black!30},
  backgroundcolor=\color{black!3},
  xleftmargin=4pt, xrightmargin=4pt,
  extendedchars=true,
  literate={á}{{\'a}}1 {é}{{\'e}}1 {í}{{\'i}}1 {ó}{{\'o}}1 {ú}{{\'u}}1
           {Á}{{\'A}}1 {É}{{\'E}}1 {Í}{{\'I}}1 {Ó}{{\'O}}1 {Ú}{{\'U}}1
           {ñ}{{\~n}}1 {Ñ}{{\~N}}1
}
 desde cada documento. Compilar con pdflatex (o latexmk -pdf).
\usepackage[utf8]{inputenc}
\usepackage[T1]{fontenc}
\usepackage{lmodern}
\usepackage[spanish,es-noshorthands,es-tabla,es-nodecimaldot]{babel}
\usepackage{amsmath,amssymb}
\usepackage[margin=2.3cm]{geometry}
\usepackage{array,booktabs,longtable}
\usepackage[normalem]{ulem}
\usepackage{xcolor}
\usepackage{enumitem}
\usepackage{listings}
\usepackage[hidelinks]{hyperref}

\setlength{\parindent}{0pt}
\setlength{\parskip}{0.5em}
\setlist{itemsep=2pt,topsep=2pt}
\renewcommand{\arraystretch}{1.25}
\Urlmuskip=0mu plus 1mu

% Símbolos frecuentes
\newcommand{\sz}{\ensuremath{\sigma^0}}
\newcommand{\gz}{\ensuremath{\gamma^0}}
\newcommand{\bz}{\ensuremath{\beta^0}}
\newcommand{\dsz}{\ensuremath{\Delta\sigma^0}}
\newcommand{\ok}{\ensuremath{\checkmark}}
\newcommand{\flecha}{\ensuremath{\rightarrow}}
\newcommand{\aprox}{\ensuremath{\approx}}
\newcommand{\por}{\ensuremath{\times}}

% Recuadro de nota (equivale a los ">" del Markdown)
\newenvironment{nota}
  {\begin{quote}\small\itshape}
  {\end{quote}}

% Bloques de código y comandos
\lstset{
  basicstyle=\ttfamily\small,
  breaklines=true,
  columns=fullflexible,
  keepspaces=true,
  frame=single,
  rulecolor=\color{black!30},
  backgroundcolor=\color{black!3},
  xleftmargin=4pt, xrightmargin=4pt,
  extendedchars=true,
  literate={á}{{\'a}}1 {é}{{\'e}}1 {í}{{\'i}}1 {ó}{{\'o}}1 {ú}{{\'u}}1
           {Á}{{\'A}}1 {É}{{\'E}}1 {Í}{{\'I}}1 {Ó}{{\'O}}1 {Ú}{{\'U}}1
           {ñ}{{\~n}}1 {Ñ}{{\~N}}1
}

\setlength{\parskip}{0.3em}

\title{TerraPulse AI\\[2pt]\large Ecuaciones del radar: medición, cambio e interpretación}
\author{Documentación técnica del proyecto}
\date{14 de noviembre de 2026}

\begin{document}
\maketitle

\begin{abstract}
El radar de apertura sintética (SAR) de NISAR no usa ecuaciones de fluidos. Su física es
electromagnética (ecuaciones de Maxwell, propagación y dispersión de microondas). Este documento
reúne las ecuaciones que intervienen en el flujo de TerraPulse AI: la ecuación del radar, la
normalización de la retrodispersión, la conversión a decibeles y el cambio entre fechas. También
incluye los modelos de dispersión que sirven para \emph{interpretar} un cambio y la relación
interferométrica que solo aplicaría si el caso de estudio fuera deformación.
\end{abstract}

\section{Ecuación del radar}

Para un radar monoestático (la misma antena emite y recibe), la potencia recibida de un blanco es
\begin{equation}
  P_r = \frac{P_t\, G^2\, \lambda^2\, \sigma}{(4\pi)^3 R^4},
  \label{eq:radar}
\end{equation}
donde
\begin{itemize}
  \item $P_t$: potencia transmitida;
  \item $G$: ganancia de la antena;
  \item $\lambda$: longitud de onda (banda L de NISAR, \emph{Frequency A}: $\lambda \approx 24$\,cm);
  \item $R$: distancia entre el radar y el blanco;
  \item $\sigma$: sección eficaz radar del blanco (m$^2$).
\end{itemize}

\section{Coeficientes de retrodispersión}

En superficies extensas (suelo, vegetación, agua) la sección eficaz se normaliza por un área de
referencia y se obtiene un coeficiente adimensional:
\begin{equation}
  \sigma^0 = \frac{\langle \sigma \rangle}{A_{\text{suelo}}} .
\end{equation}
Las tres normalizaciones habituales difieren en el área de referencia. Para terreno plano, con
ángulo de incidencia local $\theta_i$:
\begin{equation}
  \beta^0 = \frac{\sigma^0}{\sin\theta_i},
  \qquad
  \gamma^0 = \frac{\sigma^0}{\cos\theta_i}.
\end{equation}
En los productos GCOV de NISAR, $\gamma^0$ se normaliza con el área iluminada proyectada en el
plano perpendicular a la línea de vista, calculada con un modelo digital de elevación. Esto
corrige el efecto del relieve, por lo que en terreno con pendiente la relación sencilla
$\gamma^0 = \sigma^0/\cos\theta_i$ es solo una aproximación.

La capa \texttt{HHHH} de GCOV es el término diagonal de la matriz de covarianza,
$\langle |S_{HH}|^2 \rangle$: la potencia retrodispersada con emisión y recepción horizontales.

\section{Conversión a decibeles}

\begin{equation}
  \sigma^0_{\text{dB}} = 10 \log_{10} \sigma^0 .
  \label{eq:db}
\end{equation}
La misma expresión se aplica a $\gamma^0$. Si los datos fueran complejos (GSLC), primero se
calcula la potencia, $10\log_{10}|z|^2$.

\section{Cambio entre dos fechas comparables}

Para dos escenas comparables (mismo producto, trayectoria, marco, dirección, polarización,
resolución, CRID y madurez) en la misma malla:
\begin{equation}
  \Delta\sigma^0_{\text{dB}}(x,y) = \sigma^0_{\text{dB},t_2}(x,y) - \sigma^0_{\text{dB},t_1}(x,y)
  = 10\log_{10}\frac{\sigma^0_{t_2}(x,y)}{\sigma^0_{t_1}(x,y)} .
  \label{eq:delta}
\end{equation}
La diferencia en dB equivale al logaritmo de la razón de potencias.

\paragraph{Anomalía estandarizada (planeada, aún no implementada).}
Respecto a la estadística de la región o de una línea base:
\begin{equation}
  z(x,y) = \frac{\Delta\sigma^0_{\text{dB}}(x,y) - \mu}{s},
\end{equation}
donde $\mu$ y $s$ son la media y la desviación estándar de referencia. Los parámetros exactos se
documentarán en \texttt{pipeline/change.py} cuando exista.

\paragraph{Advertencia.} Un $\Delta\sigma^0$ describe un cambio en la respuesta radar; por sí solo
no identifica su causa física.

\section{Modelos de dispersión (para interpretar, no para medir)}

\subsection{Rugosidad: criterio de Rayleigh}
Una superficie se comporta como lisa (reflexión especular, poca señal de vuelta) si la desviación
estándar de su altura $h$ cumple
\begin{equation}
  h < \frac{\lambda}{8\cos\theta_i}.
\end{equation}
El agua en calma cumple este criterio y aparece oscura en la imagen. El suelo rugoso o húmedo
devuelve más señal.

\subsection{Dispersión de volumen: \emph{water cloud model}}
Pese a su nombre, no es un modelo de fluidos. Representa el dosel vegetal como una nube de
dispersores (Attema y Ulaby, 1978):
\begin{align}
  \sigma^0 &= A\, V_1 \cos\theta_i \left(1 - T^2\right) + T^2\, \sigma^0_{\text{suelo}}, \\
  T^2 &= \exp\!\left(-\frac{2 B\, V_2}{\cos\theta_i}\right),
\end{align}
donde $V_1$ y $V_2$ son descriptores de la vegetación, $A$ y $B$ son parámetros empíricos que se
ajustan con datos de campo y $T^2$ es la transmisividad de ida y vuelta a través del dosel.

\subsection{Doble rebote}
El agua bajo vegetación (vegetación inundada) forma un reflector diédrico tronco--agua que
\emph{aumenta} la señal, sobre todo en HH. Es el comportamiento opuesto al del agua abierta, y es
relevante si el fenómeno elegido fuera inundación o humedales.

\section{Interferometría (solo si el caso fuera deformación)}

La diferencia de fase entre dos adquisiciones se relaciona con el desplazamiento en la línea de
vista $\Delta R$:
\begin{equation}
  \Delta\phi = \frac{4\pi}{\lambda}\, \Delta R .
\end{equation}
El signo depende de la convención del producto (GUNW). El resultado es un desplazamiento en línea
de vista, no vertical.

\section{¿Dónde entrarían ecuaciones de fluidos?}

Solo fuera del radar, al interpretar el fenómeno: por ejemplo, un modelo hidráulico de inundación o
uno de advección--dispersión de sedimentos. Ninguno está en el plan actual. La metodología del
proyecto pide validar con datos independientes (precipitación, Sentinel-2/HLS, OPERA), no con
modelos.

\section*{Referencias}
\begin{itemize}
  \item Documentación NISAR (ASF): \url{https://nisar-docs.asf.alaska.edu/}
  \item Ulaby, F.~T. y Long, D.~G. (2014). \emph{Microwave Radar and Radiometric Remote Sensing}.
        University of Michigan Press.
  \item Attema, E.~P.~W. y Ulaby, F.~T. (1978). Vegetation modeled as a water cloud.
        \emph{Radio Science}, 13(2), 357--364.
  \item Small, D. (2011). Flattening gamma: radiometric terrain correction for SAR imagery.
        \emph{IEEE Transactions on Geoscience and Remote Sensing}, 49(8), 3081--3093.
\end{itemize}

\end{document}
